\chapter*{Abstract}
\addcontentsline{toc}{chapter}{Abstract} % Add the preface to the table of contents as a chapter

Spatial transcriptomics has emerged as a powerful technology for studying gene expression while preserving the spatial organization of cells within tissues. In cancer research, these approaches provide valuable insights into the tumour microenvironment and the interactions between tumour and immune cells. This Master's thesis focused on the analysis of Xenium spatial transcriptomics data from muscle-invasive bladder cancer (MIBC) samples obtained from the DUTRENEO clinical trial, which investigated neoadjuvant immunotherapy strategies.

The main objective of this project was to develop a robust and biologically informed workflow for the analysis of Xenium spatial transcriptomics data and to apply it to investigate mechanisms underlying differential response to immune checkpoint inhibitor (ICI) therapy. To achieve this, a review of the current state of the art in spatial transcriptomics quality control and cell annotation methodologies was first conducted. Based on this review, a comprehensive preprocessing and analysis pipeline was established, including quality control, normalization, dimensionality reduction, clustering, and cell-type annotation.

Subsequently, neighbourhood analysis was performed to characterize the spatial organization of cellular communities within the tumour microenvironment. The results revealed heterogeneous spatial patterns and distinct neighbourhood structures enriched in epithelial, fibroblast, and immune-related populations. Notably, non-responders to ICI therapy exhibited fibroblast-enriched epithelial communities consistent with an immune-excluded phenotype, suggesting that stromal barriers may limit immune-cell access to tumour regions despite the presence of inflammatory signatures. Furthermore, statistical analysis of cell-type distributions across spatial communities identified significant differences between responders and non-responders, supporting the existence of distinct tumour microenvironment architectures associated with treatment outcome.

Overall, this thesis provides a standardized workflow for Xenium spatial transcriptomics analysis and demonstrates the value of spatial transcriptomics for uncovering biological mechanisms that cannot be captured by conventional transcriptomic approaches alone. The findings provide insights into potential mechanisms of resistance to immunotherapy in MIBC and highlight the potential of spatial biomarkers to improve patient stratification and support personalized treatment strategies.

\vspace{1.8cm}

\textbf{Keywords:} biomarker, bladder cancer, cisplatin, immune checkpoint inhibitors, neoadjuvant
therapy, patient stratification, PDL1, randomized trial, multi-omics, spatial transcriptomics

% %%%%%%%%%%%%%%%%%%%%%%%%%%%%%%%%%%%%%%%%%%%%%%%%%%%%%%%%%%%%%%%%%%%%%%%%%%%%%%%%%%%%%%%%%%%%%%%
% I am of the opinion that every \LaTeX\xspace geek, at least once during 
% his life, feels the need to create his or her own class: this is what 
% happened to me and here is the result, which, however, should be seen as 
% a work still in progress. Actually, this class is not completely 
% original, but it is a blend of all the best ideas that I have found in a 
% number of guides, tutorials, blogs and tex.stackexchange.com posts. In 
% particular, the main ideas come from two sources:

% \begin{itemize}
% 	\item \href{https://3d.bk.tudelft.nl/ken/en/}{Ken Arroyo Ohori}'s 
% 	\href{https://3d.bk.tudelft.nl/ken/en/nl/ken/en/2016/04/17/a-1.5-column-layout-in-latex.html}{Doctoral 
% 	Thesis}, which served, with the author's permission, as a backbone 
% 	for the implementation of this class;
% 	\item The 
% 		\href{https://github.com/Tufte-LaTeX/tufte-latex}{Tufte-Latex 
% 			Class}, which was a model for the style.
% \end{itemize}

% The first chapter of this book is introductory and covers the most
% essential features of the class. Next, there is a bunch of chapters 
% devoted to all the commands and environments that you may use in writing 
% a book; in particular, it will be explained how to add notes, figures 
% and tables, and references. The second part deals with the page layout 
% and design, as well as additional features like coloured boxes and 
% theorem environments.

% I started writing this class as an experiment, and as such it should be 
% regarded. Since it has always been intended for my personal use, it may
% not be perfect but I find it quite satisfactory for the use I want to 
% make of it. I share this work in the hope that someone might find here 
% the inspiration for writing his or her own class.

% \begin{flushright}
% 	\textit{Federico Marotta}
% \end{flushright}
